\chapter{Realizacja}\label{ch:realizacja}

\TemplateNote{Rozdział pokazuje rysunek (makro \latexcmd{Figure}), rysunek złożony z podrysunków, listing wyjścia
terminala, listing kodu oraz długą tabelę. Pliki graficzne umieszczaj w katalogu \file{images/}.}

\section{Rysunki}\label{sec:rysunki}

Rysunek wstawia się makrem \latexcmd{Figure[szerokość]\{plik\}\{podpis\}\{etykieta\}}; plik podaje się względem
katalogu \file{images/}, a szerokość jest ułamkiem szerokości tekstu. Podpis rysunku znajduje się \emph{pod}
rysunkiem. Rysunek~\ref{fig:pojedynczy} przedstawia przykład; gdy plik nie istnieje, w jego miejscu pojawia się
ramka z nazwą brakującego pliku, więc dokument kompiluje się także bez zdjęć.

\Figure[0.6]{placeholder.pdf}{Przykładowy rysunek wstawiony makrem \latexcmd{Figure}}{fig:pojedynczy}

Kilka powiązanych obrazów łączy się w jeden rysunek środowiskiem \code{subfigure} (rysunek~\ref{fig:podrysunki}).
Odwołania do podrysunków mają postać \latexcmd{ref\{fig:podrysunek-a\}}, co daje np. rysunek~\ref{fig:podrysunek-a}.

\begin{figure}[H]
    \centering
    \begin{subfigure}[b]{0.48\linewidth}
        \centering
        \includegraphics[width=\linewidth]{placeholder.pdf}
        \caption{Pierwszy podrysunek}\label{fig:podrysunek-a}
    \end{subfigure}\hfill
    \begin{subfigure}[b]{0.48\linewidth}
        \centering
        \includegraphics[width=\linewidth]{placeholder.pdf}
        \caption{Drugi podrysunek}\label{fig:podrysunek-b}
    \end{subfigure}
    \caption{Rysunek złożony z dwóch podrysunków}\label{fig:podrysunki}
\end{figure}

\section{Listingi}\label{sec:listingi}

Domyślny styl \code{terminal} służy do wklejania wyjścia konsoli (listing~\ref{lst:terminal}). Styl jest wspólny
dla całego dokumentu, więc w pojedynczym listingu wystarczy podać podpis i etykietę.

\begin{lstlisting}[caption={Przykład wyjścia terminala (sesja w konsoli)}, label={lst:terminal}]
$ sha256sum dane.bin
e3b0c44298fc1c149afbf4c8996fb92427ae41e4649b934ca495991b7852b855  dane.bin
$ ls -l dane.bin
-rw-r--r-- 1 user user 0 paź 25 11:15 dane.bin
\end{lstlisting}

Kod źródłowy zapisuje się w stylu \code{code} z numeracją wierszy i podświetlaniem składni (listing~\ref{lst:kod}).
Polskie znaki w komentarzach i napisach działają bez dodatkowej konfiguracji. Dłuższe pliki można wczytać poleceniem
\latexcmd{lstinputlisting[style=code, language=...]\{plik\}}.

\begin{lstlisting}[style=code, language=Python, caption={Przykładowy fragment kodu w języku Python}, label={lst:kod}]
def srednia(wartosci):
    """Zwraca średnią arytmetyczną listy liczb."""
    if not wartosci:
        raise ValueError("Pusta lista wartości")
    return sum(wartosci) / len(wartosci)

print(srednia([2, 4, 6]))  # wynik: 4.0
\end{lstlisting}

\section{Długa tabela}\label{sec:dluga-tabela}

Tabele dłuższe niż jedna strona składa się środowiskiem \code{longtable}; nagłówek powtarza się wtedy
na każdej stronie (tabela~\ref{tab:dluga}). Nie używa się w nim kolumn \code{X}, tylko kolumn o stałej
szerokości (\code{L}, \code{C}, \code{R}).

{\tablefont
\begin{longtable}{@{} l L{0.5\linewidth} L{0.25\linewidth} @{}}
    \caption{Przykład tabeli rozciągającej się na kilka stron}\label{tab:dluga} \\
    \toprule
    \hcell{Lp.} & \hcell{Opis} & \hcell{Uwagi} \\
    \midrule
    \endfirsthead
    \caption[]{Przykład tabeli rozciągającej się na kilka stron (ciąg dalszy)} \\
    \toprule
    \hcell{Lp.} & \hcell{Opis} & \hcell{Uwagi} \\
    \midrule
    \endhead
    \midrule
    \multicolumn{3}{r}{\emph{ciąg dalszy na następnej stronie}} \\
    \endfoot
    \bottomrule
    \endlastfoot
    1  & Pierwszy wiersz tabeli                & uwaga \\
    2  & Drugi wiersz tabeli                   & --- \\
    3  & Trzeci wiersz tabeli                  & uwaga \\
    4  & Czwarty wiersz tabeli                 & --- \\
    5  & Piąty wiersz tabeli                   & uwaga \\
    6  & Szósty wiersz tabeli                  & --- \\
    7  & Siódmy wiersz tabeli                  & uwaga \\
    8  & Ósmy wiersz tabeli                    & --- \\
    9  & Dziewiąty wiersz tabeli               & uwaga \\
    10 & Dziesiąty wiersz tabeli               & --- \\
    11 & Jedenasty wiersz tabeli               & uwaga \\
    12 & Dwunasty wiersz tabeli                & --- \\
    13 & Trzynasty wiersz tabeli               & uwaga \\
    14 & Czternasty wiersz tabeli              & --- \\
    15 & Piętnasty wiersz tabeli               & uwaga \\
    16 & Szesnasty wiersz tabeli               & --- \\
    17 & Siedemnasty wiersz tabeli             & uwaga \\
    18 & Osiemnasty wiersz tabeli              & --- \\
    19 & Dziewiętnasty wiersz tabeli           & uwaga \\
    20 & Dwudziesty wiersz tabeli              & --- \\
    21 & Dwudziesty pierwszy wiersz tabeli     & uwaga \\
    22 & Dwudziesty drugi wiersz tabeli        & --- \\
    23 & Dwudziesty trzeci wiersz tabeli       & uwaga \\
    24 & Dwudziesty czwarty wiersz tabeli      & --- \\
    25 & Dwudziesty piąty wiersz tabeli        & uwaga \\
    26 & Dwudziesty szósty wiersz tabeli       & --- \\
    27 & Dwudziesty siódmy wiersz tabeli       & uwaga \\
    28 & Dwudziesty ósmy wiersz tabeli         & --- \\
    29 & Dwudziesty dziewiąty wiersz tabeli    & uwaga \\
    30 & Trzydziesty wiersz tabeli             & --- \\
    31 & Wiersz tabeli nr 31                   & uwaga \\
    32 & Wiersz tabeli nr 32                   & --- \\
    33 & Wiersz tabeli nr 33                   & uwaga \\
    34 & Wiersz tabeli nr 34                   & --- \\
    35 & Wiersz tabeli nr 35                   & uwaga \\
    36 & Wiersz tabeli nr 36                   & --- \\
    37 & Wiersz tabeli nr 37                   & uwaga \\
    38 & Wiersz tabeli nr 38                   & --- \\
    39 & Wiersz tabeli nr 39                   & uwaga \\
    40 & Wiersz tabeli nr 40                   & --- \\
    41 & Wiersz tabeli nr 41                   & uwaga \\
    42 & Wiersz tabeli nr 42                   & --- \\
    43 & Wiersz tabeli nr 43                   & uwaga \\
    44 & Wiersz tabeli nr 44                   & --- \\
    45 & Wiersz tabeli nr 45                   & uwaga \\
    46 & Wiersz tabeli nr 46                   & --- \\
    47 & Wiersz tabeli nr 47                   & uwaga \\
    48 & Wiersz tabeli nr 48                   & --- \\
    49 & Wiersz tabeli nr 49                   & uwaga \\
    50 & Wiersz tabeli nr 50                   & --- \\
\end{longtable}
}
